\chapter{Introduction}
\label{chap:intro}
\lhead{\emph{Introduction}}

As we’re in the era of Digital Transformation, Artificial Intelligence and Machine Learning have a real impact in the best way, using AI in cybersecurity is a huge boost in this generation. The Log Monitoring Security with AI  resolution came with a solution for monitoring logs, 
as everyday machines create a bunch of unstructured data sometimes, with the help of AI the Monitoring system will make a huge difference by solving problem of like faults positive, unusual behaviour of the system, attacks and much more security events\cite{kaur2023artificial}.

Log monitoring is a crucial aspect of cybersecurity. With organisations generating massive amounts of log data daily, identifying threats in real-time is a growing challenge. AI can help by automating the detection and resolution of anomalies, reducing the workload for security teams, and providing faster responses to potential threats,  monitoring logs in real-time and providing security for applications, infrastructure components-computers, networks, storage systems, and all other services as there are many challenges in log analysis \cite{oliner2012advances}.  It provides critical insight into system health and security by finding real-time issues, performance bottlenecks, and possible security threats. Log monitoring security collects log data, categorises filters, and analyses it for actionable insights. They can automatically trigger alerts, create reports, and even perform corrective actions when specific thresholds or anomalies are detected. It typically integrates this system with other tools such as security information and event management platforms, DevOps pipes, or incident management systems that ensure the smooth running of operations\cite{pratap2022real}.

 

The idea behind this is to come up with a system that will monitor servers security in real-time using Artificial Inteligent, to make log monitoring a friendly interface where you don't need great IT Skills to see what is going on the system. The Log monitoring will help by  being able to see in clear text and  get hint how to solve a problem on the system. the rise of AI really changes how we do things now by implementing AI for security monitoring resolution give details about how to solve a problem or how to provide a solution to the engineer or the person who is managing system. This is one of the features the Log monitoring system security will provide for users. \cite{turnbull2005understanding}





 

 


\section{Motivation}

The reason for me to design this system the idea came from my work placement in Trend Micro. I was working as a Cybersecurity Technical Support Engineer, and my focus was on customer satisfaction by solving problems they were facing in Cloud One Workload Security \cite{trendmicro_business}. 



To solve customer problems the first information we ask is collecting the log of the endpoint causing a problem and the next step to do is to manually checking every single file to detect problem, as humain sometime you can miss some event due to size of log received from each server and we had a tool that was helping us for the first troubleshoot the way it was working you have to collect logs from the affected server and upload it to the tools and the tolls will check every single file to find error that can tell us what the customer has on the system. based on that  and dealing with many different cases I was inspired to come up with the idea of AI to help for Log monitoring will be a great boost for everyone in IT and Cybersecurty\cite{stanciu2013importance}.

In my time at Trend Micro I had a case where a customer wanted to create a rule to find all failed authentication attempts, but he was not able to do it because the Cloud One agent could not allow him to meet his requirements. After work, I went home thinking about it, and I wrote a Python script that could meet the exact requirements, it worked fine, and I started doing a search about Log Inspection, which was one of the Cloud One modules. By doing this search, I found out how useful it is to have a Log Monitoring Security with AI resolution that can allow you to create any rules for any application you would run on your server, in my case, I’ll be more focused on the security and networking solve different issue people have for a better troubleshooting a solution to run applications safely and easily\cite{hassani2018studying}. As a Technical Support Engineer, troubleshooting is the most important part of solving problems most of the time; engineers check logs manually to find problems sometimes; there are many different tools they use, like Splunk an example, and troubleshooting is a long process, to facilitate things for everyone sides customers and Engineer side that why I believe Log Monitoring system with real-time events management can make a big change and make life easier for everyone. Since security is crucial, I was inspired to add a security feature that can lead to quick threat research on the systems by monitoring some security events that are needed with the help of AI to provide a first-step solution\cite{chainourov2017log}. 


%Why is it important to do a project on this topic? This should cover your key motivation for this. For example an excellent student from 2016 noticed a large number of homeless sleeping rough in Cork and was motivated to develop a system that load balanced the homeless shelters to try to accommodate the maximum number of homeless. This section can include the personal pronoun but the rest of the report should be third person passive, this is the case with most technical reports! For example here it is fine to say "... I decided to develop and app to help ...".

\section{Contribution}






%Enumerate ck themtand how to trahe main contributions. Here try to zoom out, to talk from the perspective of a Computer Science graduate. In other words, imagine you are talking to a job panel, and you want to show your computer science skills by enumerating how they are reflected in your project work. A good guide here is to look back over the modules you have covered as an undergrad from 2/3rd year, how many tools and techniques from these modules do you have in the project and to what extent? How have you advanced beyond the module content? Do you have anything new?

The foundation of Log Monitoring Security with AI resolution is based on a combination of the whole program of IT Management and Cybersecurity, nearly every single module from the first year to the final year has contributed to the realisation of this project. 


Operating System module helped me to understand how each OS works, This project will be more focus on Unix/Linux operating system\cite{zeng2016computer}. 

Linux Administration, by learning how to administrate Linux servers it gives a good understanding how to troubleshoot some problems that can put a server in panic mode, Intrusion Detection system it will be use as part of the project and the system is build around that too, that will make things easier to see who is attacking, and do some emergency defense on a server\cite{8474119}. 

The Log Monitoring Security need to have a web console, from this module Server-side web development it helps to contacted the server with the log by sending requested as we have feature of request sedurity report the server side of the console will help to make some API Calls, so  the user can be able to see what going on they have to use a console, to build the console Server-side Web Development helps by building a web application that can handle all the information will be collected and displayed them so the user can be able to see how the data is monitor.

Database Design: To be able to collect and process all the logs, they need to be placed somewhere safe and well structure so they can be filtered and make things easier to be accessible, a good design will be so important to be implemented, because the data collected will be unstructured and will be transform by a structure one. the module helps to have a good understanding how LMS can deal with data and logs collected from the system.

Programming the system used to build on programming langue like python in fields like \textbf{AI/ML, cybersecurity, big data analytics, and automation}.





\section{Structure of This Document}
% notice how I cross referenced the chapters through using the \label tag --> LaTeX is VERY similar to HTML and other mark up languages so you should see nothing new here!
This document contains 7 Chapter, from the Introduction to Conclusion each chapter give you an insist into the project and system will be implement later on, you'll get a clear understating of Log Monitoring Security with AI resolution, The historic, problem facing in today world, from this you'll understand how important is to using the system for a better IT Management and Security

The first chapter of this project is where I introduce you on this this project explaining the basic and clear idea of this project, the next section is the Motivation where I elaborate about what motivate me to work on this project, how everything started to be inspired and the desire of doing searched about the topic and Finally the contribution where you can see how this project contribute not just in the Computer Science and real business out there.

The Second Chapter is where you can have a clear understanding about my project where I explain about area of computer science where this project landed, I also explain about real business , enterprise sector this system can be use.

The Third Chapter you can find what problem this project about to solved, the objective to archive, functional requirements and Non Functional requirements.

Chapter fourth is all about the implementation Approach, you can find the Architecture of the system, how it's going to be build, Risk Assessments where there are some potential risk, Methodology, Implementation Plan Schedule, Evaluation and Prototype of the system.

The Chapter Five is about the Conclusion and Future Work, we'll see some problem faced durent the project. 


%This section is quite formulaic. Briefly describe the structure of this document, enumerating what does each chapter and section stands for. For instance in this work in Chapter \ref{chap:background} the guidance in structuring the literature review is given. Chapter \ref{chap:problem} describes the main requirements for the problem definition and so on ...